\chapter{Appendix A - Loan Application DCR Graph}
\label{app:loanapp_graph}

This appendix presents the process model discovered from the Loan
Application event log introduced in Section~\ref{sec:loan_dataset},
using Disco's Fuzzy Miner algorithm. The model captures the
observed control-flow structure in terms of activity frequencies and
transitions, with node size reflecting activity frequency and arc
thickness representing transition frequency. The resulting model is
shown in Figure~\ref{app:loanapp_disco}.

Table~\ref{tab:loanapp_relations} complements this view by providing
a full explicit enumeration of the DCR relation set used in the
formal model presented in the main text.

\begin{figure}[H]
    \centering
    \includegraphics[width=0.9\textwidth]{Pictures/LoanApp-Disco.png}
    \caption{Process model discovered by Disco (Fuzzy Miner) from the
             Kirchdorfer et al.\ log. Node size encodes activity
             frequency; arc thickness encodes transition frequency.}
    \label{app:loanapp_disco}
\end{figure}

\begin{longtable}{p{5.5cm} p{5.5cm} p{2cm}}
\caption{Complete relation set of the Loan Application \ac{DCR} graph:
23 conditions, 11 responses, and 27 exclusions (61 relations in total).
There are no include, milestone, or co-response relations.}
\label{tab:loanapp_relations}\\

\toprule
\textbf{Source} & \textbf{Target} & \textbf{Relation} \\
\midrule
\endfirsthead

\multicolumn{3}{c}{\tablename\ \thetable\ -- continued from previous page} \\
\toprule
\textbf{Source} & \textbf{Target} & \textbf{Relation} \\
\midrule
\endhead

\midrule
\multicolumn{3}{r}{Continued on next page} \\
\endfoot

\bottomrule
\endlastfoot

\multicolumn{3}{l}{\textit{Conditions (23)}} \\
\midrule
Check application form completeness & Appraise property & Condition \\
Check application form completeness & Check credit history & Condition \\
Check application form completeness & AML check & Condition \\
Appraise property & Assess loan risk & Condition \\
Check credit history & Assess loan risk & Condition \\
AML check & Assess loan risk & Condition \\
Applicant completes form & Assess loan risk & Condition \\
Assess loan risk & Design loan offer & Condition \\
Check application form completeness & Design loan offer & Condition \\
Design loan offer & Approve loan offer & Condition \\
Appraise property & Approve loan offer & Condition \\
Check credit history & Approve loan offer & Condition \\
AML check & Approve loan offer & Condition \\
Approve loan offer & Approve application & Condition \\
Check application form completeness & Approve application & Condition \\
Assess loan risk & Approve application & Condition \\
Assess loan risk & Reject application & Condition \\
Return application to applicant & Reject application & Condition \\
Approve loan offer & Cancel application & Condition \\
Check application form completeness & Cancel application & Condition \\
Assess loan risk & Cancel application & Condition \\
Check application form completeness & Return application to applicant & Condition \\
Return application to applicant & Applicant completes form & Condition \\

\addlinespace
\multicolumn{3}{l}{\textit{Responses (11)}} \\
\midrule
Check application form completeness & Appraise property & Response \\
Check application form completeness & Check credit history & Response \\
Check application form completeness & AML check & Response \\
Appraise property & Assess loan risk & Response \\
Check credit history & Assess loan risk & Response \\
AML check & Assess loan risk & Response \\
Design loan offer & Approve loan offer & Response \\
Return application to applicant & Applicant completes form & Response \\
Applicant completes form & Check application form completeness & Response \\
Applicant completes form & Assess loan risk & Response \\
Applicant completes form & Reject application & Response \\

\addlinespace
\multicolumn{3}{l}{\textit{Exclusions (27)}} \\
\midrule
Appraise property & Return application to applicant & Exclusion \\
Appraise property & Applicant completes form & Exclusion \\
Check credit history & Return application to applicant & Exclusion \\
Check credit history & Applicant completes form & Exclusion \\
AML check & Return application to applicant & Exclusion \\
AML check & Applicant completes form & Exclusion \\
Assess loan risk & Check application form completeness & Exclusion \\
Design loan offer & Reject application & Exclusion \\
Design loan offer & Appraise property & Exclusion \\
Design loan offer & Check credit history & Exclusion \\
Design loan offer & AML check & Exclusion \\
Approve loan offer & Assess loan risk & Exclusion \\
Approve application & Cancel application & Exclusion \\
Approve application & Design loan offer & Exclusion \\
Reject application & Design loan offer & Exclusion \\
Reject application & Approve loan offer & Exclusion \\
Reject application & Approve application & Exclusion \\
Reject application & Cancel application & Exclusion \\
Reject application & Check credit history & Exclusion \\
Reject application & Appraise property & Exclusion \\
Reject application & AML check & Exclusion \\
Cancel application & Approve application & Exclusion \\
Cancel application & Design loan offer & Exclusion \\
Return application to applicant & Design loan offer & Exclusion \\
Return application to applicant & Approve loan offer & Exclusion \\
Return application to applicant & Approve application & Exclusion \\
Return application to applicant & Cancel application & Exclusion \\
\end{longtable}


\chapter{Appendix B - Role Extension XML Schema}
\label{app:roleoptions_xml}

Section~\ref{sec:roles} introduces the \texttt{roleOptions} extension
to the \ac{DCR}-JS XML schema, used to attach per-role cost and
duration multipliers to an event. Listing~\ref{lst:roleoptions_xml}
shows the exact schema used in the Loan Application experiment
(Table~\ref{tab:role_multipliers}).

\begin{lstlisting}[language=XML, caption={\texttt{roleOptions} XML extension, Loan Application Junior/Expert profiles.}, label={lst:roleoptions_xml}]
<roleOptions>
  <role name="Junior"
        costMultiplier="0.7"
        durationMultiplier="2.0"/>
  <role name="Expert"
        costMultiplier="2.0"
        durationMultiplier="0.5"/>
</roleOptions>
\end{lstlisting}


\chapter{Appendix C - Step Reward Computation for the Expense Report Pareto Traces}
\label{app:expense_step_reward}

The step rewards reported in the \textbf{Step reward} column of
Table~\ref{tab:pareto_traces} follow directly from the normalised
reward of Equation~\ref{eq:reward_total} (Section~\ref{sec:reward_normalised}),
computed under $\alpha=0.0$, $\beta=1.0$, with $d_{\max}=7{,}200$\,min
(the duration of \texttt{Approve}, the longest event in this benchmark):

\medskip
\noindent
\begin{tabular}{@{}ll@{}}
\texttt{FillOut} (intermediate): &
  $1 - 1.0 \times \frac{1{,}500}{7{,}200}
  = 1 - 0.208 = +0.79$ \\[4pt]
\texttt{Approve} (intermediate): &
  $1 - 1.0 \times \frac{7{,}200}{7{,}200}
  = 1 - 1.0 = \phantom{+}0.00$ \\[4pt]
\texttt{PayOut} (accepting): &
  $1 - 1.0 \times \frac{1{,}800}{7{,}200} + 100
  = 0.75 + 100 = +100.75$ \\[4pt]
\texttt{Withdraw} (accepting): &
  $1 - 1.0 \times \frac{1{,}200}{7{,}200} + 100
  = 0.833 + 100 = +100.83$ \\
\end{tabular}

\medskip
\noindent
The step rewards illustrate the normalised reward function operating
as intended: \texttt{FillOutExpenseReport} incurs a duration penalty
in both traces, while the accepting event delivers the terminal bonus
regardless of the path taken. P2 reaches acceptance in two steps at a
higher monetary cost, while P1 takes three steps at lower cost but
longer duration --- precisely the trade-off the Pareto front captures.

\section{Per-Decile Reward Convergence}
\label{app:expense_reward_deciles}

Section~\ref{sec:expense_baseline} states that the mean per-episode
reward rises sharply during the first decile of training and then
stays within a narrow band for the remainder, placing the final
20\,\% evaluation window (Section~\ref{sec:training_pipeline}) well
inside the converged regime. Table~\ref{tab:expense_reward_deciles}
reports the underlying mean reward per training decile, for all six
weight configurations, supporting that claim with the exact values.

\begin{table}[H]
\centering
\footnotesize
\caption{Mean per-episode reward by training decile, Expense Report
benchmark, all six weight configurations. Decile 1 is the first
10\,\% of training; decile 10 is the last (the final two deciles
together form the 20\,\% evaluation window used throughout
Chapter~\ref{chap:results}).}
\label{tab:expense_reward_deciles}
{\renewcommand{\arraystretch}{1.4}
\setlength{\tabcolsep}{4pt}
\begin{tabular}{lrrrrrrrrrr}
\toprule
\textbf{$(\alpha,\beta)$} & \textbf{D1} & \textbf{D2} & \textbf{D3} & \textbf{D4} & \textbf{D5} & \textbf{D6} & \textbf{D7} & \textbf{D8} & \textbf{D9} & \textbf{D10} \\
\midrule
$(0.0, 0.0)$ & 89.2 & 98.1 & 99.1 & 99.6 & 99.9 & 100.3 & 100.9 & 100.5 & 100.8 & 100.7 \\
$(1.0, 0.0)$ & 89.0 & 98.0 & 99.0 & 99.4 & 99.8 & 100.6 & 100.2 & 100.9 & 98.4  & 97.3  \\
$(0.0, 1.0)$ & 89.6 & 98.2 & 97.6 & 96.4 & 98.8 & 98.9  & 99.2  & 99.2  & 96.9  & 96.9  \\
$(0.5, 0.5)$ & 89.9 & 98.3 & 99.1 & 97.4 & 99.3 & 99.4  & 99.6  & 99.4  & 98.0  & 99.4  \\
$(2.0, 0.5)$ & 90.0 & 98.3 & 99.1 & 99.3 & 99.1 & 99.4  & 99.3  & 96.1  & 99.1  & 98.3  \\
$(0.5, 2.0)$ & 88.4 & 97.9 & 96.8 & 94.5 & 98.6 & 98.8  & 97.4  & 96.8  & 98.7  & 98.6  \\
\bottomrule
\end{tabular}}
\end{table}


\chapter{Appendix D - Loan Application Results}
\label{app:loanapp_results}

This appendix collects three complementary results for the
role-conditioned Loan Application experiment of
Section~\ref{sec:results_loanapp}: additional training diagnostics,
the complete recovered Pareto-optimal set, and the role choice along a
representative trace under each Expert capacity budget.

\section{Training Diagnostics}
\label{app:loanapp_diagnostics}

These figures are omitted from Section~\ref{sec:results_loanapp}
because compliance and shielding behaviour are analysed in
Chapter~\ref{chap:internal_validation}.

\begin{figure}[H]
  \centering
  \includegraphics[width=\textwidth]{Pictures/loanapp_reward_convergence.png}
  \caption{Reward convergence on the role-conditioned Loan Application
  benchmark across scalarisation weight pairs.}
  \label{fig:app_loanapp_reward}
\end{figure}

\begin{figure}[H]
  \centering
  \includegraphics[width=\textwidth]{Pictures/loanapp_accept_rate.png}
  \caption{Accept rate over training on the role-conditioned Loan
  Application benchmark.}
  \label{fig:app_loanapp_accept}
\end{figure}

\begin{figure}[H]
  \centering
  \includegraphics[width=\textwidth]{Pictures/loanapp_shielding_activity.png}
  \caption{Shielding activity over training, measured as the percentage
  of steps that were illegal action attempts.}
  \label{fig:app_loanapp_shielding}
\end{figure}

\begin{figure}[H]
  \centering
  \includegraphics[width=\textwidth]{Pictures/loanapp_clean_episodes.png}
  \caption{Percentage of accepting episodes with zero illegal action
  attempts on the role-conditioned Loan Application benchmark.}
  \label{fig:app_loanapp_clean}
\end{figure}

\section{Complete Pareto-Optimal Set}
\label{app:loanapp_pareto_points}

Table~\ref{tab:app_loanapp_pareto_points} reports the complete
Pareto-optimal set recovered by the full framework on the
role-conditioned Loan Application benchmark.

\begin{table}[H]
\centering
\renewcommand{\arraystretch}{1.5}
\setlength{\tabcolsep}{14pt}
\caption{Complete 26-point Pareto-optimal set recovered on the
role-conditioned Loan Application benchmark (unlimited capacity, seed
$s=1$), pooled across the five non-baseline weight configurations,
last 20\,\% of training. Underlies the summary in
Table~\ref{tab:loanapp_pareto_summary} and
Figure~\ref{fig:loanapp_pareto_balanced}.}
\label{tab:app_loanapp_pareto_points}
\begin{tabular}{rrrl}
\toprule
\textbf{\#} & \textbf{Cost} & \textbf{Duration} & \textbf{Configuration} \\
\midrule
1  & 1\,092.0 & 1\,632.0 & $\alpha=1.0,\ \beta=0.0$ \\
2  & 1\,190.8 & 1\,549.5 & $\alpha=0.5,\ \beta=0.5$ \\
3  & 1\,245.4 & 1\,506.0 & $\alpha=0.5,\ \beta=0.5$ \\
4  & 1\,344.2 & 1\,423.5 & $\alpha=0.5,\ \beta=0.5$ \\
5  & 1\,505.4 & 1\,369.5 & $\alpha=0.5,\ \beta=0.5$ \\
6  & 1\,604.2 & 1\,287.0 & $\alpha=0.5,\ \beta=0.5$ \\
7  & 1\,658.8 & 1\,243.5 & $\alpha=0.5,\ \beta=0.5$ \\
8  & 1\,757.6 & 1\,161.0 & $\alpha=0.5,\ \beta=0.5$ \\
9  & 1\,907.1 & 1\,140.0 & $\alpha=0.5,\ \beta=0.5$ \\
10 & 1\,946.1 & 1\,119.0 & $\alpha=0.5,\ \beta=2.0$ \\
11 & 2\,005.9 & 1\,057.5 & $\alpha=0.5,\ \beta=0.5$ \\
12 & 2\,044.9 & 1\,036.5 & $\alpha=0.5,\ \beta=0.5$ \\
13 & 2\,167.1 & 1\,003.5 & $\alpha=0.5,\ \beta=0.5$ \\
14 & 2\,206.1 &   982.5  & $\alpha=0.5,\ \beta=0.5$ \\
15 & 2\,265.9 &   921.0  & $\alpha=0.0,\ \beta=1.0$ \\
16 & 2\,304.9 &   900.0  & $\alpha=0.5,\ \beta=0.5$ \\
17 & 2\,320.5 &   877.5  & $\alpha=0.5,\ \beta=0.5$ \\
18 & 2\,359.5 &   856.5  & $\alpha=0.5,\ \beta=2.0$ \\
19 & 2\,419.3 &   795.0  & $\alpha=0.5,\ \beta=2.0$ \\
20 & 2\,458.3 &   774.0  & $\alpha=0.0,\ \beta=1.0$ \\
21 & 2\,607.8 &   753.0  & $\alpha=0.5,\ \beta=0.5$ \\
22 & 2\,706.6 &   670.5  & $\alpha=0.0,\ \beta=1.0$ \\
23 & 2\,867.8 &   616.5  & $\alpha=0.0,\ \beta=1.0$ \\
24 & 2\,966.6 &   534.0  & $\alpha=0.0,\ \beta=1.0$ \\
25 & 3\,021.2 &   490.5  & $\alpha=0.0,\ \beta=1.0$ \\
26 & 3\,120.0 &   408.0  & $\alpha=0.0,\ \beta=1.0$ \\
\bottomrule
\end{tabular}
\end{table}

\section{Role Choice Along a Single Trace, by Expert Capacity Budget}
\label{app:loanapp_capacity_trace}

Figure~\ref{fig:loanapp_capacity_allocation} (Section~\ref{sec:loanapp_capacity})
reports the Expert-budget allocation pattern pooled over many episodes.
Figure~\ref{fig:loanapp_capacity_trace} shows the same mechanism at the
level of a single trace, comparing the most-balanced accepting episode
at each budget $k$ (the point closest to the ideal cost--duration
corner after min-max normalisation, the same criterion used for
Table~\ref{tab:loanapp_pareto_summary}). The Expert budget is a single
pool shared across the whole episode, depleted by one each time it is
chosen, regardless of which event consumes it.

\begin{figure}[H]
  \centering
  \includegraphics[width=0.85\textwidth]{Pictures/LoanApp_capacity_trace.png}
  \caption{Role choice along the most-balanced accepting trace, by Expert capacity budget $k$.}
  \label{fig:loanapp_capacity_trace}
\end{figure}

At $k=1$ and $k=2$ the trace spends the entire budget on the
highest-sensitivity events identified in Section~\ref{sec:loanapp_roles}
(Check application form completeness, then Check credit history),
leaving the rest of the trace on Junior. At $k=\infty$, the event order
itself changes (Appraise property and Check credit history swap
position), showing that removing the budget constraint affects not
only role choice but also the sequencing the policy prefers.

At $k=3$, the most-balanced trace uses only 2 of the 3 available
Expert assignments. This is not under-utilisation: across the
recovered front, episodes using all 3 Expert assignments reach lower
minimum duration (616 vs.\ 879) but their cost and duration ranges
overlap substantially with those using only 2, since duration also
depends on event order. The balanced point's distance-to-ideal
criterion selects the configuration with the best joint trade-off, not
the one that maximises Expert usage; the agent treats the scarce
budget as a resource to allocate selectively rather than a quota to
exhaust.
